\section{Attack Surface Analysis}


\subsection{Real-World Applications}

\textbf{Mozilla Firefox.} Firefox uses libpng to decode PNG images
rendered in web content. PNG data is processed incrementally as it is
downloaded over HTTP, and passed to libpng through its progressive
decoding interface (\texttt{png\_process\_data()}). This resembles the
progressive png decoding strategy used in our fuzzing harness. Because any
website can supply arbitrary image data via \texttt{<img>} tags,
background images, or favicons, the library is exposed to untrusted
input in a fully remote, browser-driven context.

\textbf{GIMP.} The GNU Image Manipulation Program relies on libpng
to decode and encode PNG files when opening or exporting images.
When a user loads an image from an untrusted source (e.g., email
attachments or downloaded assets), GIMP parses PNG metadata,
including iTXt chunks, using libpng’s standard read interfaces.
This exposes the same parsing logic exercised in our test harness
to user-controlled files in a desktop application context.

\subsection{Concrete Attack Scenario}

An attacker hosts or distributes a crafted PNG image containing a
malformed iTXt chunk. The \textit{Length} field is set to match the
declared keyword size, but the data omits a NUL terminator within the
allocated buffer.

When a victim application such as a web browser or image viewer
processes the image, libpng parses the file and enters the iTXt
decoding routine in \texttt{png\_push\_read\_iTXt()}. During this
process, pointer-based traversal is used to locate NUL separators
between metadata fields. Because no terminator exists within bounds, the parser continues
reading beyond the allocated heap buffer. This results in a small
out-of-bounds read during string processing (via \texttt{strlen()}
or equivalent traversal logic), leading to memory access outside the
intended allocation.

In practice, this can result in a controlled crash (denial of
service) or limited disclosure of adjacent heap memory, depending on
the application environment and memory layout. The issue requires no user interaction
beyond opening or rendering the image.

\subsection{Code Paths Not Exercised by the Harness}

\textbf{Gap 1: Encoder (write) path.}
The fuzzing harness exercises only read-side PNG processing through
\texttt{png\_process\_data()} and does not invoke any write-side APIs
such as \texttt{png\_write\_info()} or \texttt{png\_write\_iTXt()}.
As a result, the encoder code in \texttt{pngwutil.c} is not covered.

This matters in applications such as GIMP, where PNG metadata may be
read from a file and later written back when exporting an image.

\textbf{Gap 2: Optional image transform paths.}
The harness does not enable certain image transformation routines,
such as RGB-to-grayscale conversion (\texttt{png\_do\_rgb\_to\_gray()}
in \texttt{pngtrans.c}). These paths may be triggered in real-world
applications under specific rendering conditions (e.g., grayscale
display modes or image filters).

This gap implies that certain execution paths are not explored, and
therefore potential vulnerabilities in these components would not be
triggered by the current fuzzing campaign.